\documentclass{iiair}
\usepackage[utf8]{inputenc}

\usepackage{csquotes}
\usepackage{hyperref} 
\usepackage{tikz}

\usepackage[numbers]{natbib}

% Delete this package if unnecessary
\usepackage{lipsum} 

% ------- IIAIR Editor ---------
\iiairReceived{11 Jul 2024}  
\iiairAccepted{25 Aug 2024}
\iiairPublished{1 Dec 2024}
\iiairFirstPage{1}
\iiairVolume{1}
\iiairIssue{1}
\iiairYear{2024}
\iiairdoi{https://doi.org/10.2478/iiair}
\iiaireissn{XXXX-XXXX}
% ------------------------------

\title{This is the Title Format for Innovative Informatics and Artificial Intelligence Research (IIAIR)}
\shorttitle{Short version of title}

\author[1]{First author\,\orcidlink{0000-0000-0000-0000}\,}
\author[1]{Second author\,\orcidlink{0000-0000-0000-0000}}
\author[2]{Third author}

\affil[1]{Department, University, City,
    Country}
\affil[2]{Department, University, City}
  
\iiairauthorcorr{First Author (\href{mailto:author1@itk.ac.id}{author1@itk.ac.id})}
\shortauthor{Author et al.}
\shorttitle{IIAIR LaTeX Template}


 
\bibliographystyle{IEEEtran}

\begin{document}

% -------------- Highlighted Contributions ------------
\thispagestyle{empty}
\begin{center}
    \huge\textbf{Highlight of Contributions}
\end{center}

\bigskip

% Add your contributions here
\begin{minipage}{\textwidth}
\large \textbf{Auhor 1, Author 2, Author 3}
\begin{itemize}
\large
    \item contribution 1
    \item contribution 2
    \item contribution 3
\end{itemize}
\end{minipage}
% ----------------------------------------------------


\maketitle

\begin{abstract}
  These are guidelines for preparing papers for the \emph{Innovative Informatics and Artificial Intelligence Research (IIAIR)}. The journal accepts original research
  papers, survey, and tutorial articles within the broad field of informatics and/or artificial intelligence. Use this document as a template if you are using
  \LaTeX; otherwise, use this document as an
  instruction set. The paper size is A4 (21 × 29.7\,cm).  

  \keywords{artificial intelligence; machine learning; deep learning; informatics}
\end{abstract}

\section{Introduction}


This document is a template for \LaTeX. 
Please use the electronic version of this document as a
template when you produce your manuscript for submission to the
\emph{Innovative Informatics and Artificial Intelligence Research (IIAIR)}. The paper size is A4 (21 × 29.7~cm).

These are the examples of subsection

\subsection{Sub section 1}

This is a subsection of introduction
\subsubsection{Sub section 1.1}
This is a subsection of sub-section 1
\section{Related Works}

This section is intended for your literature review and recent discussion.

\section{Methods}

In the methods section, please describe all necessary information about your methods.

This is an example of figure template in IIAIR LaTeX format. Please use bold style text in a caption as shown in Figure \ref{fig:transformer}.

Figures may be in color or black and white and must be of such quality
that they produce clear and sharp printouts on an ordinary (color)
laser printer.

\begin{figure}[t]
  \centering
  \includegraphics[width=0.7\columnwidth]{transformer.png}
  \caption{\textbf{The transformer network model \cite{attention_transformer2017}}}
  \label{fig:transformer}
\end{figure}

\section{Results and Discussios}
This section is intended for the results and your discussion.

The table format can be seen in Table \ref{tbl:example}. The table's style is by adding a top and bottom line in the title for each column and a bottom line at the end of the content in the table. The presentation of the table in a single-column format can be achieved by appending an asterisk (*) to the table command.

\begin{table}[t]
\caption{\textbf{Example of table format}}
\label{tbl:example}
\centering
\begin{tabular}{ccc}
\hline
head1    & head2    & head3    \\ \hline
content1 & content2 & content3 \\
content1 & content2 & content3 \\ \hline
\end{tabular}
\end{table}




\section{Conclusion}
Conclusion section
\section{Acknowledgments}
Acknowledgements section if necessary


% Editor Check
\section{Delete this Section: Lorem}
\lipsum[2-5]


\newpage

\bibliography{references}
\end{document}
